Theorem~\ref{thm:circuit-bound} establishes that, given a program
that terminates with a specified complexity bound, one can extract a
concrete quantum circuit exhibiting the same space complexity. This result
is particularly significant, as it enables the characterization of
 quantum polynomial time, as shown in Theorem~\ref{thm:fbqp-sound} and
 Theorem~\ref{thm:fbqp-complete}. However, the burden of proving termination of
 a \hyrql{} program currently rests entirely on the programmer.
This naturally raises the following question: to what
extent is the termination of \hyrql{} programs semi automatable?

We argue that such automatization is feasible by leveraging techniques
from Term Rewrite Systems (TRS)\footnote{More specifically, a higher-order version
called Simply-Typed TRS~\cite{Yam01}, but we do not enter this level of details
here.}, as they
are equipped with a rich toolbox for analyzing both termination and
computational complexity. Approaches to termination analysis include
non-exhaustively recursive path
orderings~\cite{Der82,Der87}, dependency pairs~\cite{AG00}, and
size-change termination~\cite{LJBA01}. Similarly, several methods
address complexity analysis, e.g.,
interpretation methods~\cite{BMM11,MP09,BDL16}, and polynomial path
orders~\cite{Moser09,AM09}.

Although a full formal treatment is left to future work, we
present an illustrative example supporting our claim. Specifically, we
illustrate how a \hyrql{} program can be efficiently translated into a (quantum) TRS,
how the reduction behaviour of the original program is faithfully
simulated by the resulting rewriting system while preserving
complexity bounds, and how termination and complexity results
established for the TRS can be transferred back to the corresponding
\hyrql{} program.


For example, the standard NOT gate is represented by the quantum TRS whose set
of rewrite rules is $ \{\mathtt{NOT}\,\zket \to \oket, \mathtt{NOT}\,\oket \to
\zket\}$. The application of the $\mathtt{NOT}$ function symbol to $\zket$
rewrites to $\oket$. Similarly, the Hadamard gate $H$ can be defined using a
superposition of terms $ \{\mathtt{H}\,\zket \to \frac{1}{\sqrt 2} \cdot \zket
+ \frac{1}{\sqrt 2} \cdot \oket, \mathtt{H}\,\oket \to  \frac{1}{\sqrt 2} \cdot
\zket - \frac{1}{\sqrt 2} \cdot \oket\}$. The evaluation of a quantum term will
follow a parallel strategy, similar to \hyrql{} operational semantics
(Table~\ref{tab:reduction}), which reduces in parallel the distinct branches of
a superposition, when possible (otherwise, the reduction will have an
exponential blowup).

We claim that we can automatize the translation of any term of \hyrql{} to a
corresponding TRS. This translation works as follows:

\begin{itemize}
  \item Produce a set of rules, inductively on the syntax of the term;
  \item Yield a function symbol $\mtf$ for the initial term $t$;
  \item When encountering $\qcasedefault x$ (same for \textmatch{}), merge the 
    two sets of rules generated for $t_0$ and $t_1$, while keeping track of
    which path was taken.
\end{itemize}

Coming back on Example~\ref{ex:bqwalk}, such an algorithm would generate the
following rewrite rules:
\[
  \left\{
    \begin{aligned}
      \mathtt{Repeat}\,0 &\to \nil   &\mathtt{BQWalk}\,\zket\,m &\to \zket ::
      \mathtt{Repeat}\,m \\
      \mathtt{Repeat}\,S(m) &\to \zket:: \mathtt{Repeat}\,m  &
      \mathtt{BQWalk}\,\oket\,0 &\to \oket :: \nil \\
      & &     \mathtt{BQWalk}\,\oket\,S(m) &\to \oket ::
      \mathtt{BQWalk}\,(\ket -)\,m
    \end{aligned}
  \right\}
\]

One can remark that this set of rules captures completely
the behaviour of $\mathtt{bqwalk}$, and they both compute the same function;
furthermore, one reduction step for a given rewrite rule will correspond to
$1 \leq k \leq \size{\mathtt{bqwalk}}$ reduction steps in \hyrql{}.
Therefore, a certificate on this TRS yields directly a certificate on the
original term. 

Now, it is possible to apply the tools we mentioned to the obtained TRS. In
particular, we can use recursive path orderings~\cite{Der82} to show
termination, coupled with quasi-interpretations~\cite{BMM11} to get that this
TRS terminates in time $\mcal O(n)$ for any input of size $n$. While such
properties are not defined for quantum TRS, i.e.,\ with superpositions and
parallel reduction, they can be adapted to fit the quantum setting; a formal
adaptation is left for future work. Altogether with
Theorem~\ref{thm:circuit-bound}, we can conclude that
$\mathtt{bqwalk}$ can be compiled to a family of quantum circuits, whose size
is polynomial in the size of the input.

Using the same ideas, we claim that we can describe a semantic-preserving
translation algorithm from \hyrql{} to TRS, with no blowup on the termination
complexity in either direction. Such work, along with a proper development of
TRS in the quantum setting and an adaptation of existing techniques, would give
a semi-automatized way to prove the feasibility/tractability of terms in
\hyrql{}.

While this section was focused on termination and time complexity, some of the
mentioned techniques also provide results to bound the space complexity of a
TRS program~\cite{BMM11}. Given a good definition of quantum space, this would
make it possible to characterize quantum space complexities, and obtain bounds
on the depth or width of the circuit.
